% sec-07: UC San Diego Moores Cancer Center.  No figure by design: application
% 10 already carries the intake sequence, the authority diagram, the feasibility
% grid, and the approval dependency.
\section{UC San Diego Moores Cancer Center}

\subsection{The Positioning Constraint}

UC San Diego Moores Cancer Center is named in all ten applications as the
intended partner of choice, approached at the feasibility stage only. It is not
described in any of them, or in this paper, as a partner, sponsor, trial site,
or endorser, and it will not be so described without written authorization. No
agreement of any kind exists. The one application addressed to the institution
asks for a 45-minute meeting.

The constraint is not a disclaimer bolted on at the end. It shapes what the
applications can claim: a funder reading application 01 is told the site is
sought, not secured, and the budget's first-year gate is written as \textit{site
executed} rather than as \textit{site confirmed}. An application that overstated
the relationship would fail its own go or no-go condition in year one.

This section carries no figure. Application 10 already contains the intake
sequence, the role-authority diagram, the feasibility question grid, and the
approval dependency, and drawing any of them again here would be duplication
rather than a new perspective. The master prompt's requirement that every figure
be a unique perspective is a constraint on adding figures, not a quota to fill.

\subsection{Three Corrections the Site Research Forced}

The most useful thing the site research produced was not a contact list. It was
a list of things the applications were getting wrong.

\begin{apptable}
\begin{tabularx}{\textwidth}{>{\raggedright\arraybackslash}p{3.8cm} >{\raggedright\arraybackslash}X >{\raggedright\arraybackslash}p{3.4cm}}
\headrow \thead{Claim as first drafted} & \thead{Correction} & \thead{Where it now appears}\\
\midrule
Daraxonrasib is first-in-human & It is investigational and already in Phase 3 evaluation. The supportable claim is the first prospective clinical evaluation of the integrated surgical and advisory workflow, subject to FDA and institutional confirmation & A positioning note in nine applications; a pre-send checklist line in all ten \\
An eight-arm robotic Whipple & No site can assess feasibility or regulatory status without manufacturer, model, cleared configuration, arm count, instruments, and software version. The configuration is specified at the site agreement & Application 10 §1; the positioning note in the other nine \\
Drug access is available & No drug supply, letter of authorization, or regulatory cross-reference is in place with the agent's developer & Application 06 pre-send instruction; application 10 §1 and back matter \\
\end{tabularx}
\end{apptable}
\tabcap{Three corrections the site research forced, and where each now appears.
Nine of the ten applications carried the first error until the PART I audit
pass, which is recorded here rather than quietly fixed.}

The first correction is worth dwelling on because it is the kind of error that
survives review. \textit{First-in-human} was true of the integrated workflow and
false of the drug, and in a title the two readings are indistinguishable. A
reviewer at a drug-development office would have caught it immediately and
discounted everything after it. The correction cost two sentences; not making it
would have cost the application.

\subsection{Two Routes, Sequenced Differently}

The repository carries research on two candidate sites, and they are sequenced
in opposite directions on purpose.

\begin{apptable}
\begin{tabularx}{\textwidth}{>{\raggedright\arraybackslash}p{3.0cm} >{\raggedright\arraybackslash}X >{\raggedright\arraybackslash}X}
\headrow \thead{} & \thead{UC San Diego} & \thead{Scripps}\\
\midrule
Entry point & Pancreatic surgery and GI medical oncology & AI validation and digital-trial methods \\
First ask & A feasibility meeting on the trial & An independent review of the AI portfolio \\
The trial appears at & Step 3 of 12 & Step 9 of 13 \\
Risk if the route fails & The trial has no site & The AI work has no independent reviewer \\
\end{tabularx}
\end{apptable}
\tabcap{The two site routes on four properties. They are sequenced in opposite
directions because they are asking different institutions for different things,
not because one is a fallback for the other.}

The UC San Diego route is trial-first because Moores Cancer Center has the
pancreatic surgical volume and the trial operations the study needs, and the
question worth asking there is whether the concept is clinically supportable at
all. The Scripps route is AI-first because the value it can add is independent
evaluation of the verification methods, and asking a methods group to host a
surgical trial would waste the introduction.

Application 03 uses the contrast directly. Its pathway figure shows the branch
taken if the first site agreement does not execute, and that branch is a second
qualified site rather than the end of the programme, because the performer, not
the site, holds the regulatory package.

\subsection{What a Productive Early Outcome Looks Like}

The site research sets a deliberately modest bar, and the applications inherit
it. A productive outcome is an invitation to present, the identification of a
clinical principal investigator, acceptance into an institutional feasibility
process, a preliminary budget and accrual estimate, a technical review of the
robotic and AI workflow, a defined regulatory and contracting pathway, or
support for a grant application. None of those is a trial, and none of them is
described as one.

A patient-facing trial proceeds only after the site, the manufacturer, the FDA
pathway, the funding, the sponsor responsibilities, the contracts, the
insurance, the institutional review board approval, and the operational
resources are all documented. That sentence is in the site research, in every
application's back matter, and in this paper, and it is the same sentence in all
three places.
